\chapter{Marco Teórico}

En esta sección se presentan los fundamentos conceptuales y técnicos que sustentan el desarrollo de esta investigación. Se definen los principios teóricos, modelos y enfoques metodológicos que permiten comprender el contexto en el que se sitúa el estudio, así como las nociones esenciales relacionadas con el análisis de logs de seguridad, los modelos de lenguaje de gran escala y el razonamiento agentivo aplicado a la ciberseguridad. Estos elementos proporcionan el marco conceptual necesario para sustentar tanto la revisión de la literatura como el diseño de la propuesta desarrollada.

%─────────────────────────────────────────────────────────────────
\section{Logs de seguridad}

Los logs de seguridad constituyen registros textuales generados de forma continua por sistemas operativos, aplicaciones, dispositivos de red y servicios en la nube, que documentan eventos relevantes ocurridos en la infraestructura tecnológica de una organización. Cada entrada de log contiene información estructurada o semiestructurada sobre la naturaleza del evento, su origen, su severidad y el momento en que ocurrió.

Desde una perspectiva técnica, un log puede representarse como una tupla de la forma:

\begin{equation}
    l_i = (\tau_i,\ s_i,\ \rho_i,\ m_i)
\end{equation}

\noindent donde $\tau_i \in \mathbb{R}$ es el \textit{timestamp} del evento, $s_i$ identifica la fuente generadora (host, servicio o dispositivo), $\rho_i$ corresponde al nivel de severidad (DEBUG, INFO, WARNING, ERROR, CRITICAL) y $m_i \in \Sigma^*$ es el mensaje de texto no estructurado que describe el evento.

El valor de los logs radica en que permiten reconstruir la secuencia de eventos ocurridos en un sistema, facilitando la detección de comportamientos anómalos, la auditoría de accesos y la investigación forense de incidentes. Sin embargo, en entornos de producción modernos, una sola infraestructura puede generar cientos de miles de eventos por hora provenientes de fuentes heterogéneas, lo que hace inviable su revisión manual \cite{palma2025llm}. Esta escalabilidad es precisamente uno de los desafíos centrales que motiva el uso de técnicas automatizadas de análisis.

%─────────────────────────────────────────────────────────────────
\section{Parsing de logs}

El \textit{log parsing} es el proceso de transformar mensajes de log no estructurados en representaciones estructuradas, extrayendo plantillas (\textit{templates}) fijas y parámetros variables. Este paso es fundamental para cualquier sistema de análisis automatizado, ya que permite identificar el tipo de evento independientemente de los valores específicos que contiene cada instancia \cite{ma2024llmparser}.

El proceso de \textit{parsing} transforma cada entrada $l_i$ en una tupla estructurada:

\begin{equation}
    \hat{l}_i = (\tau_i,\ s_i,\ \rho_i,\ T_i,\ P_i)
\end{equation}

\noindent donde $T_i$ es la plantilla extraída que representa la estructura fija del mensaje y $P_i = \{p_1, \ldots, p_k\}$ es el conjunto de parámetros variables identificados, tales como direcciones IP, nombres de usuario o identificadores de proceso.

Los enfoques tradicionales de \textit{parsing}, basados en expresiones regulares o árboles de prefijos como Drain, presentan limitaciones frente a la diversidad y evolución de los formatos de log en sistemas modernos. En este contexto, los modelos de lenguaje de gran escala han emergido como alternativas más flexibles, capaces de generalizar a nuevos formatos sin reconfiguración manual \cite{ma2024llmparser, liu2025librelog}. Estudios recientes reportan que modelos como LLaMA y Flan-T5 alcanzan precisiones de \textit{parsing} superiores al 96\% en colecciones de 16 sistemas de código abierto \cite{ma2024llmparser}.

%─────────────────────────────────────────────────────────────────
\section{Detección de anomalías en logs}

La detección de anomalías en logs consiste en identificar eventos o secuencias de eventos que se desvían significativamente del comportamiento normal esperado del sistema. Estos eventos anómalos pueden indicar fallos de software, intrusiones, movimientos laterales de un atacante o exfiltración de datos \cite{almodovar2024logfit}.

Formalmente, dado un conjunto de \textit{embeddings} $\{\phi(\hat{l}_i)\}_{i=1}^n$, se define una función de puntuación de anomalía:

\begin{equation}
    \text{AnomalyScore}(\phi(\hat{l}_i)) \in [0, 1]
\end{equation}

\noindent donde valores superiores a un umbral $\theta$ determinado por el modelo de detección indican un evento anómalo. Los enfoques actuales se dividen en tres categorías principales: supervisados, que requieren datos etiquetados; no supervisados, que aprenden la distribución normal sin etiquetas; e híbridos, que combinan ambos paradigmas \cite{liu2024logcem}.

Los métodos basados en modelos de lenguaje han demostrado ventajas significativas sobre los enfoques tradicionales, especialmente en escenarios donde los logs contienen variabilidad semántica que no puede capturarse con expresiones regulares o modelos de recuento de eventos. \cite{almodovar2024logfit} proponen LogFiT, un modelo basado en BERT fine-tuneado para detección de anomalías mediante predicción enmascarada de tokens, logrando superar a los baselines en los datasets HDFS, BGL y Thunderbird sin requerir datos etiquetados.

%─────────────────────────────────────────────────────────────────
\section{Sistemas SIEM y Centros de Operaciones de Seguridad}

Un \textit{Security Information and Event Management} (SIEM) es una plataforma tecnológica que centraliza, normaliza y correlaciona eventos de seguridad provenientes de múltiples fuentes de la infraestructura tecnológica de una organización, con el propósito de detectar amenazas y facilitar la respuesta a incidentes en tiempo real \cite{kim2026isiem}. Los SIEMs modernos integran capacidades de agregación de logs, correlación de eventos mediante reglas, generación de alertas y almacenamiento histórico para análisis forense.

Un \textit{Security Operations Center} (SOC) es la unidad organizacional que opera el SIEM y coordina la respuesta a incidentes de seguridad. Los analistas del SOC son responsables de monitorear las alertas generadas, investigar los eventos sospechosos, contener las amenazas identificadas y documentar los incidentes \cite{srinivas2025soc}.

Sin embargo, los SIEMs tradicionales basados en reglas estáticas presentan limitaciones estructurales bien documentadas en la literatura. \cite{oduola2025wazuh} evalúan la plataforma Wazuh contra el framework MITRE ATT\&CK y reportan una tasa de detección global del 85\%, con brechas significativas en tácticas de Recolección y Exfiltración, lo que evidencia la incapacidad de los sistemas basados en reglas para correlacionar cadenas de ataque multi-fase. Adicionalmente, el fenómeno de \textit{alert fatigue}, provocado por el alto volumen de falsos positivos, reduce significativamente la efectividad de los equipos de seguridad \cite{palma2025llm}.


\subsection{Vulnerabilidades conocidas: CVE y NVD}

Un \textit{Common Vulnerability and Exposure} (CVE) es un identificador
estándar asignado públicamente a una vulnerabilidad de seguridad conocida
en software o hardware, que permite referenciarla de forma unívoca entre
distintas herramientas, bases de datos y equipos de respuesta. Cada
entrada CVE incluye una descripción de la vulnerabilidad, los productos
afectados y referencias a recursos adicionales de mitigación.

La \textit{National Vulnerability Database} (NVD) es el repositorio
oficial del Instituto Nacional de Estándares y Tecnología de Estados
Unidos (NIST) que enriquece cada entrada CVE con métricas de severidad
estandarizadas mediante el sistema CVSS (\textit{Common Vulnerability
Scoring System}), rangos de versiones afectadas y referencias a parches
disponibles. La NVD expone una API pública que permite consultar
vulnerabilidades por palabra clave, producto o identificador CVE,
facilitando la integración con sistemas automatizados de análisis de
seguridad. En el agente propuesto, esta API se utiliza para enriquecer
el diagnóstico con las vulnerabilidades conocidas asociadas a los
servicios identificados en los logs analizados.

%─────────────────────────────────────────────────────────────────
\section{Modelos de Lenguaje de Gran Escala}

Los Modelos de Lenguaje de Gran Escala (\textit{Large Language Models}, LLMs) son modelos de inteligencia artificial basados en la arquitectura \textit{Transformer}, entrenados sobre grandes corpus de texto para aprender representaciones estadísticas del lenguaje natural y de código fuente. Su capacidad para comprender y generar texto de forma coherente los ha convertido en herramientas versátiles para una amplia gama de tareas de procesamiento de lenguaje natural \cite{karlsen2023llm4sec}.

En el contexto del análisis de logs de seguridad, los LLMs ofrecen ventajas significativas sobre los enfoques tradicionales: pueden capturar relaciones semánticas complejas entre eventos, adaptarse a nuevos formatos de log sin reconfiguración manual y generar explicaciones interpretables para los analistas de seguridad \cite{palma2025llm}. \cite{karlsen2023llm4sec} evalúan cinco arquitecturas de LLM (BERT, RoBERTa, DistilRoBERTa, GPT-2 y GPT-Neo) para análisis de logs de seguridad, reportando que el modelo DistilRoBERTa fine-tuneado alcanza un F1-Score promedio de 0.998 en seis datasets de logs de sistemas y aplicaciones web.

\subsection{Fine-tuning}

El \textit{fine-tuning} es el proceso de adaptar un LLM preentrenado a un dominio específico mediante entrenamiento adicional sobre datos del dominio objetivo. Este proceso permite al modelo especializar sus representaciones internas para las particularidades léxicas y semánticas del dominio, mejorando significativamente su rendimiento en tareas especializadas \cite{wang2024finellm}. En el análisis de logs, el fine-tuning permite que el modelo aprenda los patrones específicos de los formatos de log de un sistema particular, mejorando tanto la precisión del \textit{parsing} como la detección de anomalías.

\subsection{Prompt Engineering}

El \textit{prompt engineering} consiste en el diseño sistemático de instrucciones y contextos (\textit{prompts}) para guiar el comportamiento de un LLM hacia tareas específicas, sin necesidad de reentrenamiento. Esta técnica es especialmente valiosa en escenarios \textit{zero-shot} o \textit{few-shot}, donde no se dispone de suficientes datos etiquetados para el fine-tuning \cite{liu2024logprompt}. \cite{liu2024logprompt} demuestran que estrategias avanzadas de \textit{prompting} mejoran el rendimiento de los LLMs en análisis de logs hasta en un 380.7\% comparado con prompts simples, logrando superar a enfoques entrenados con miles de ejemplos sin requerir datos de dominio.

%─────────────────────────────────────────────────────────────────
\section{Recuperación Aumentada por Generación (RAG)}

La \textit{Retrieval-Augmented Generation} (RAG) es una técnica que combina la recuperación de información relevante desde una base de conocimiento externa con la generación de respuestas de un LLM. El proceso consiste en: (1) transformar la consulta en un vector semántico mediante un modelo de embeddings; (2) recuperar los documentos más relevantes de la base de conocimiento mediante búsqueda por similitud vectorial; (3) incluir los documentos recuperados como contexto en el \textit{prompt} del LLM para guiar la generación de la respuesta \cite{blefari2026cyberrag}.

En el contexto de la ciberseguridad, RAG permite a los LLMs acceder a información actualizada sobre tácticas y técnicas de ataque sin necesidad de reentrenamiento. \cite{singh2026cloudtrail} demuestran que integrar MITRE ATT\&CK como fuente de conocimiento recuperable mejora la precisión en la detección y mapeo de amenazas desde un 46\% hasta un 78\% en logs de AWS CloudTrail. \cite{blefari2026cyberrag} proponen CyberRAG, un sistema agentivo RAG para clasificación y reporte de ciberataques que integra grafos de conocimiento con recuperación vectorial, demostrando la viabilidad del enfoque en escenarios de análisis de amenazas reales.

%─────────────────────────────────────────────────────────────────
\section{MITRE ATT\&CK}

MITRE ATT\&CK (\textit{Adversarial Tactics, Techniques, and Common Knowledge}) es un framework de conocimiento estructurado que documenta las tácticas, técnicas y procedimientos (TTPs) utilizados por actores maliciosos en ataques reales. Organizado en una matriz de tácticas (columnas) y técnicas (filas), proporciona un vocabulario común para describir el comportamiento de los atacantes en diferentes fases de un ataque, desde el reconocimiento inicial hasta el impacto final \cite{oduola2025wazuh}.

El framework es ampliamente utilizado en ciberseguridad defensiva para mapear alertas de seguridad a tácticas y técnicas conocidas, facilitando la comprensión del contexto de un ataque y la priorización de la respuesta. En el contexto de esta investigación, MITRE ATT\&CK v15 se utiliza como base de conocimiento para el componente RAG del agente, permitiendo que el sistema correlacione eventos de log con tácticas y técnicas conocidas y genere narrativas investigativas estructuradas \cite{singh2026cloudtrail, liu2026ctithinker}.

El conocimiento estructurado de MITRE ATT\&CK se distribuye
oficialmente en formato \textbf{STIX~2.1} (\textit{Structured Threat
Information eXpression}), un estándar abierto para representar e
intercambiar inteligencia de amenazas mediante objetos JSON-LD
interoperables. Cada táctica, técnica, subtécnica, grupo de amenaza
y herramienta del framework se codifica como un objeto STIX con
identificador único, relaciones tipadas hacia otros objetos y metadatos
de procedencia. Esta representación permite poblar bases de datos
vectoriales directamente desde el repositorio CTI oficial, habilitando
la recuperación semántica de conocimiento adversarial estructurado
en tiempo de inferencia sin necesidad de reentrenamiento del modelo.

%─────────────────────────────────────────────────────────────────
%─────────────────────────────────────────────────────────────────
\section{Agentes LLM y patrones de razonamiento agentivo}
%─────────────────────────────────────────────────────────────────

Un agente LLM es un sistema de inteligencia artificial que utiliza
un modelo de lenguaje de gran escala como motor de razonamiento
para planificar y ejecutar acciones de forma autónoma, con el
objetivo de resolver tareas complejas que requieren múltiples
pasos~\cite{hassan2026elisar}. A diferencia de los pipelines LLM
convencionales, que procesan la entrada de forma lineal y producen
una única salida estática, los agentes implementan ciclos iterativos
de razonamiento y acción que les permiten adaptarse dinámicamente
a nueva información, consultar fuentes externas y refinar sus
conclusiones progresivamente.

Existen diferentes patrones arquitectónicos para implementar
agentes LLM. Los dos más relevantes para este trabajo son ReAct
y Plan-and-Execute.

El patrón \textbf{ReAct} (\textit{Reasoning + Acting}) intercala
pasos de razonamiento explícito con acciones concretas sobre el
entorno, decidiendo el siguiente paso únicamente en función de lo
observado hasta ese momento. Este enfoque es efectivo para tareas
de corta duración donde el contexto es limitado, pero presenta
limitaciones ante problemas complejos: al no disponer de una visión
global del problema, el agente puede tomar decisiones subóptimas en
cada paso, fenómeno denominado \textit{miopía de razonamiento}~\cite{rigaki2023cage}.

El patrón \textbf{Plan-and-Execute}, adoptado en la propuesta de
este seminario, resuelve esta limitación separando explícitamente la
fase de planificación de la fase de ejecución. Un módulo
\textit{Planificador} construye primero un plan completo de pasos
sobre un resumen estructurado del problema antes de ejecutar ninguna
acción; luego, un módulo \textit{Ejecutor} lleva a cabo cada paso
del plan invocando las herramientas necesarias mediante
\textit{function calling} nativo con el modelo de lenguaje. Esta
separación permite al agente mantener coherencia global a lo largo
de análisis complejos y facilita la incorporación de mecanismos de
replanificación cuando la evidencia disponible no es
suficiente~\cite{lee2025logrespagent}.

%─────────────────────────────────────────────────────────────────
\subsection{Sistemas multiagente}
%─────────────────────────────────────────────────────────────────

Un sistema multiagente (SMA) es una arquitectura en la que múltiples
agentes especializados colaboran para resolver una tarea que excede
las capacidades de un agente único. Cada agente opera sobre un
subconjunto acotado del problema, con un rol definido, un conjunto
de herramientas propio y un \textit{prompt} de sistema que delimita
su función dentro del flujo general~\cite{hmimou2025multiagent}.
Esta especialización de roles reduce la ambigüedad de las
instrucciones, mejora la consistencia de las salidas y permite
que cada agente optimice su razonamiento para una subtarea
específica, en contraste con un agente generalista que aborda
todas las responsabilidades bajo el mismo contexto.

 
%─────────────────────────────────────────────────────────────────
\section{LangGraph}
%─────────────────────────────────────────────────────────────────
 
LangGraph es un framework de Python para construir aplicaciones
agentivas con LLMs mediante la definición explícita de grafos de
estados. A diferencia de los frameworks de agentes basados en
cadenas lineales, LangGraph representa el flujo de un agente como
una máquina de estados formal (\texttt{StateGraph}), donde cada
nodo corresponde a una función Python (un agente LLM o una
herramienta determinística) y las transiciones entre nodos pueden
ser directas o condicionales~\cite{blefari2026cyberrag}.
 
El estado compartido del grafo se define como un
\texttt{TypedDict}, un diccionario tipado de Python que fluye entre
todos los nodos garantizando que cada componente opere sobre
información actualizada. Las aristas condicionales, denominadas
\textit{routers}, evalúan el estado actual y determinan
dinámicamente el nodo destino en función de los valores de sus
campos. LangGraph incorpora además un mecanismo de persistencia de
estado (\texttt{MemorySaver}) que permite registrar el historial
completo de ejecución, habilitando la trazabilidad del razonamiento
del agente.
 
Esta arquitectura es particularmente adecuada para implementar el
patrón Plan-and-Execute en entornos de análisis de seguridad, donde
el flujo de razonamiento depende de decisiones condicionales: si el
diagnóstico requiere contexto histórico, si un servicio afectado
tiene CVEs conocidas, o si los comandos generados superan la
validación de seguridad. La capacidad de LangGraph para representar
estos flujos condicionales de forma explícita y reproducible lo
posiciona como el framework de orquestación más adecuado para el
agente propuesto.


%─────────────────────────────────────────────────────────────────
\section{Bases de datos vectoriales}
%─────────────────────────────────────────────────────────────────
 
Una base de datos vectorial es un sistema de almacenamiento
especializado en la indexación y recuperación eficiente de
vectores de alta dimensión mediante métricas de similitud, como la
similitud coseno o la distancia euclidiana. A diferencia de las
bases de datos relacionales, que operan sobre valores exactos, las
bases de datos vectoriales permiten realizar búsquedas por
\textit{similitud semántica}: dado un vector de consulta, el
sistema retorna los $k$ vectores más cercanos en el espacio
vectorial~\cite{blefari2026cyberrag}.
 
En el contexto de sistemas agentivos con RAG, las bases de datos
vectoriales actúan como memoria de largo plazo del agente:
almacenan representaciones vectoriales de documentos o patrones
previos y los recuperan cuando la consulta actual presenta
similitud semántica con ellos. Esta capacidad es fundamental para
implementar \textit{few-shot} dinámico, donde los ejemplos
inyectados al modelo se seleccionan automáticamente en función de
su relevancia para el problema actual.
 
\textbf{ChromaDB} es una base de datos vectorial de código abierto
optimizada para aplicaciones de IA generativa. Soporta operación
en modo embebido (sin servidor adicional) o en modo servidor
persistente, y se integra de forma nativa con modelos de embeddings
de Sentence Transformers. Los documentos se organizan en
\textit{colecciones} independientes, cada una con su propio índice
vectorial. La recuperación se realiza mediante la función
\texttt{collection.query()}, que retorna los $k$ documentos más
similares al vector de consulta con un umbral mínimo de similitud
configurable. En la propuesta de este seminario, ChromaDB
administra tres colecciones: \texttt{attack\_patterns} para patrones
de ataque confirmados, \texttt{error\_memory} para correcciones
históricas y \texttt{host\_history} para el historial de
diagnósticos por host.

%─────────────────────────────────────────────────────────────────
\section{Embeddings semánticos}

Los embeddings semánticos son representaciones vectoriales densas de texto en un espacio de alta dimensión, donde textos con significado similar se mapean a vectores cercanos en dicho espacio. A diferencia de representaciones dispersas como TF-IDF o \textit{bag-of-words}, los embeddings semánticos capturan relaciones contextuales y semánticas entre términos, permitiendo medir la similitud conceptual entre textos \cite{liu2024logcem}.

En el análisis de logs, los embeddings permiten representar cada evento normalizado como un vector denso:

\begin{equation}
    \phi(\hat{l}_i) = \text{SentenceTransformer}(T_i \oplus P_i) \in \mathbb{R}^{768}
\end{equation}

\noindent donde $\oplus$ denota la concatenación textual del template $T_i$ y sus parámetros $P_i$, y $d = 768$ corresponde a la dimensión del modelo \texttt{all-mpnet-base-v2}. Esta representación vectorial permite aplicar algoritmos de detección de anomalías sobre el espacio semántico de los eventos, identificando logs que se desvían del comportamiento normal no solo en términos léxicos sino también conceptuales \cite{chen2025gelog}.